Este trabalho apresenta a modelagem, implementação e análise de um sistema fuzzy Mamdani para o controle de ventilação em ambientes internos. O sistema utiliza como variáveis de entrada a temperatura, a umidade relativa e o número de pessoas, determinando o nível de ventilação adequado por meio de funções de pertinência triangulares e uma base de regras linguísticas fundamentada em conforto térmico. A metodologia contempla a definição dos universos de discurso, a construção das funções de pertinência, a inferência fuzzy e a defuzzificação dos resultados. Os experimentos realizados demonstram que o sistema é capaz de tomar decisões qualitativas e adaptativas, respondendo de forma robusta às diferentes condições ambientais e evidenciando a aplicabilidade da lógica fuzzy em problemas de automação e conforto ambiental.